% GENERATED FILE. Edit canonical lessons, then run npm run generate:books.
% canonical-source-hash: fnv1a64:bd13f6a8aa920eae
% canonical-lessons: LA-C284-debes, LA-C284-explico, LA-R284-advice-and-reasons

\chapter{What You Should Do, and Why}
\label{ch:la-what-you-should-do-and-why}
\hlchaptermodality{284}

\begin{chapteropening}
\textbf{By the end of this chapter:} \emph{I can give advice with dēbēs, and explain why with quia and quod.}

Complete LA-R284-advice-and-reasons: I can give advice with dēbēs, and explain why with quia and quod.
\end{chapteropening}

\section[dēbēs, nōn dēbēs; dēbēbās]{dēbēs — what you should do}

\label{lesson:LA-C284-debes}



\begin{quote}
You own \textbf{dēbeō}, \textquotedblleft{}I have to, I must; I owe\textquotedblright{}, and \textbf{quiēscō}, \textquotedblleft{}to rest\textquotedblright{}. Say \textquotedblleft{}I have to work\textquotedblright{}. (\emph{Labōrāre dēbeō.})

Turn the same verb on a friend, and it becomes advice.
\end{quote}

\begin{grammarlens}[title={Grammar Lens: dēbēs + infinitive}]
\begin{quote}
\textbf{Aeger es; quiēscere dēbēs.} — You are ill; you should rest. \textbf{Nōn dēbēs tantum labōrāre.} — You shouldn't work so much.
\end{quote}

\emph{Dēbēs} is the second person, \textquotedblleft{}you must\textquotedblright{}, and Latin uses it for what English softens into \textquotedblleft{}you should\textquotedblright{}. The infinitive carries the advice, and \emph{nōn} in front turns it round.

\textbf{oportet}, \textquotedblleft{}it is right, one ought\textquotedblright{}, says the same thing without pointing at anyone, and it takes the accusative and infinitive you met with \emph{putō}:

\begin{quote}
\textbf{Oportet tē quiēscere.} — You ought to rest.
\end{quote}
\end{grammarlens}

\begin{grammarlens}[title={Grammar Lens: dēbēbās, you should have}]
English moves the past onto the infinitive: \textquotedblleft{}you should \textbf{have come}\textquotedblright{}. Latin moves it onto \emph{dēbeō} instead, and keeps the plain infinitive:

\begin{quote}
\textbf{Herī venīre dēbēbās.} — You should have come yesterday.
\end{quote}

\emph{Dēbēbās} is the imperfect, \textquotedblleft{}you were bound to\textquotedblright{}, and the sentence carries the regret: the coming was owed, and it did not happen. The perfect does the same work, \emph{venīre dēbuistī}, and so does \emph{oportet} in the perfect: \emph{oportuit tē venīre}.
\end{grammarlens}

\begin{cousinweb}
\textbf{dēbeō} is contracted from \emph{dē}, \textquotedblleft{}away from\textquotedblright{}, and \emph{habeō}, \textquotedblleft{}to have\textquotedblright{} (Lewis \& Short), as its own lesson said: to have something that belongs to someone else. A debt, \emph{dēbitum}, is what you hold and must give back; advice with \emph{dēbēs} treats the right action as something owed.
\end{cousinweb}

\subsection*{Your turn}
Say these aloud:

\begin{itemize}
\raggedright
  \item \emph{Aeger es; quiēscere dēbēs.}
  \item \emph{Nōn dēbēs tantum labōrāre.}
  \item \emph{Oportet tē quiēscere.}
  \item \emph{Herī venīre dēbēbās.}
  \item \emph{Fessus es; ergō quiēscere dēbēs.}
\end{itemize}

\begin{tcolorbox}[breakable,colback=teal!4,colframe=teal!35!black,title={Before you move on}]
\textquotedblleft{}You should rest\textquotedblright{}? (\emph{\textbf{Quiēscere dēbēs.}}) \textquotedblleft{}You shouldn't work so much\textquotedblright{}? (\emph{\textbf{Nōn dēbēs tantum labōrāre.}}) \textquotedblleft{}You should have come yesterday\textquotedblright{}? (\emph{\textbf{Herī venīre dēbēbās.}}) Which verb carries the past in that last one? (\emph{\textbf{Dēbēbās}}, not the infinitive.)
\end{tcolorbox}
\section[cūr? — quia …, quod …, nam …]{cūr? — quia, quod, nam: explaining}

\label{lesson:LA-C284-explico}



\begin{quote}
Ask \textquotedblleft{}why?\textquotedblright{}. (\emph{Cūr?}) Answer with \textquotedblleft{}because\textquotedblright{}. (\emph{Quia.})

You also own \textbf{explicō}, \textquotedblleft{}to explain\textquotedblright{}, and \textbf{causa}, \textquotedblleft{}a reason\textquotedblright{}.
\end{quote}

\begin{grammarlens}[title={Grammar Lens: the answer to cūr}]
\begin{quote}
\textbf{Cūr herī nōn vēnistī?} — Why didn't you come yesterday? \textbf{Nōn vēnī, quia aeger eram.} — I didn't come, because I was ill.
\end{quote}

Both pasts: the not-coming is an event, \emph{nōn vēnī}; the illness is the scene behind it, \emph{aeger eram}. After \emph{quia} the verb is an ordinary one, in whatever tense the facts need.

\textbf{nam}, \textquotedblleft{}for\textquotedblright{}, hangs the reason on afterwards, as a second sentence: \emph{Nōn vēnī. Nam aeger eram.} And \emph{explicō} names what you are doing: \textbf{Causam explicō}, I am explaining the reason.
\end{grammarlens}

\subsection*{You'll want to know: quod, because}
The four lines you wrote with joiners asked you for \emph{quia}, because \textbf{quod} has two jobs. Now you meet its second one as a reader:

\begin{quote}
\textbf{Gaudeō quod vēnistī.} — I am glad that you came, glad because you came.
\end{quote}

After a verb of feeling, \emph{gaudeō}, \emph{doleō}, \emph{quod} gives the cause. Where a noun stands before it and \emph{quod} points back to it, it is usually \textquotedblleft{}which\textquotedblright{}; where no noun waits for it, it is usually \textquotedblleft{}because\textquotedblright{} or \textquotedblleft{}that\textquotedblright{}. Writing, \emph{quia} is still the safe choice; reading, expect both.

\begin{cousinweb}
\textbf{explicō} is \emph{ex}, \textquotedblleft{}out\textquotedblright{}, with \emph{plicō}, \textquotedblleft{}to fold\textquotedblright{} (EtymDB 2.1), as its own lesson said: to explain is to unfold. A Roman book was a \emph{volūmen}, a roll, and Cicero writes \emph{volūmen explicāre} for opening one (Lewis \& Short). An explanation opens out what was rolled up inside the question.
\end{cousinweb}

\subsection*{Your turn}
Say these aloud:

\begin{itemize}
\raggedright
  \item \emph{Cūr herī nōn vēnistī?} — then answer: \emph{Nōn vēnī, quia aeger eram.}
  \item the same reason with \emph{nam} — \emph{Nōn vēnī. Nam aeger eram.}
  \item \emph{Gaudeō quod vēnistī.}
  \item \emph{Quiēscere dēbēs, quia fessus es.}
  \item a reason of your own for something you did not do yesterday
\end{itemize}

\begin{tcolorbox}[breakable,colback=teal!4,colframe=teal!35!black,title={Before you move on}]
\textquotedblleft{}Why didn't you come?\textquotedblright{} (\emph{\textbf{Cūr nōn vēnistī?}}) \textquotedblleft{}Because I was ill\textquotedblright{}? (\emph{\textbf{Quia aeger eram.}}) Which word adds the reason as a second sentence? (\emph{\textbf{Nam.}}) In \emph{gaudeō quod vēnistī}, what does \emph{quod} mean? (\textbf{Because}, or \textquotedblleft{}that\textquotedblright{}: no noun waits for it.)
\end{tcolorbox}
\section[(review)]{Review — what you should do, and why}

\label{lesson:LA-R284-advice-and-reasons}



\begin{quote}
Say \textquotedblleft{}you should rest\textquotedblright{}. (\emph{Quiēscere dēbēs.}) Then \textquotedblleft{}because you are tired\textquotedblright{}. (\emph{Quia fessus es.})
\end{quote}

\begin{grammarlens}[title={What you've built}]
Two chapters, and you can give reasons. An opinion: \emph{putō} with the accusative and infinitive, \emph{meā sententiā}, \emph{mihi vidētur}. Advice: \emph{dēbēs} and an infinitive, and \emph{dēbēbās} for what should have happened. A reason: \emph{quia}, \emph{nam}, and the \emph{quod} that follows a feeling. A conclusion: \emph{ergō}.

\begin{quote}
\emph{Meā sententiā, quiēscere dēbēs, quia aeger es. Herī labōrāre nōn dēbēbās.} — In my opinion you should rest, because you are ill. You shouldn't have worked yesterday.
\end{quote}

Four pieces in two sentences: the opinion, the advice, the reason, and the regret.
\end{grammarlens}

\subsection*{Your turn}
Say these aloud:

\begin{itemize}
\raggedright
  \item \emph{Meā sententiā, quiēscere dēbēs, quia aeger es.}
  \item \emph{Herī labōrāre nōn dēbēbās.}
  \item \emph{Cūr nōn vēnistī? — Nōn vēnī. Nam aeger eram.}
  \item \emph{Putō tē fessum esse; ergō quiēscere dēbēs.}
  \item advice of your own to a friend, with its reason
\end{itemize}

\begin{tcolorbox}[breakable,colback=teal!4,colframe=teal!35!black,title={Before you move on}]
\textquotedblleft{}You shouldn't work so much\textquotedblright{}? (\emph{\textbf{Nōn dēbēs tantum labōrāre.}}) \textquotedblleft{}You should have come\textquotedblright{}? (\emph{\textbf{Venīre dēbēbās.}}) \textquotedblleft{}I am glad that you came\textquotedblright{}? (\emph{\textbf{Gaudeō quod vēnistī.}}) \textquotedblleft{}In my opinion\textquotedblright{}? (\emph{\textbf{Meā sententiā.}})
\end{tcolorbox}
